\chapter{Direcciones y UDT}\label{ap:direcciones}

\section{Los módulos de la estación}
Según el manual de la CP-Lab (apartado 11.2.1) y la configuración de la estación Out, estas direcciones son iguales en todas las estaciones:

\begin{center}
\small
\begin{tabularx}{\textwidth}{>{\bfseries}p{4.3cm}p{2.7cm}X}
\toprule
Parte de la máquina & Entradas / Salidas & Contenido \\
\midrule
Aplicación (1 byte) & \cod{\%E0} / \cod{\%A0} & Señales propias de cada estación (cambia de una a otra). \\
Módulo básico & \cod{\%E1} / \cod{\%A1} & Botonera, selector, seta, sensores de cinta; luces, cinta, stopper. \\
Lector RFID (IO-Link) & \cod{\%E10} / \cod{\%A10} & Estado, error, dirección y 28 bytes de datos. \\
Módulo del stopper (IO-Link) & \cod{\%E42} & Sensores inductivos, encoder, «stopper abajo». \\
\bottomrule
\end{tabularx}
\end{center}

\section{Los UDT bit a bit}
\subsection*{\texttt{UDT\_Entradas\_Basico}}
Tag \texttt{Entradas\_Basico}, desde \texttt{\%I1.0}.\par\smallskip
\begin{tabular}{>{\ttfamily\small}p{5.2cm}>{\ttfamily\small}p{4.2cm}>{\ttfamily\small}p{2.4cm}}
\toprule
\textnormal{\textbf{Miembro}} & \textnormal{\textbf{Tipo}} & \textnormal{\textbf{Dirección}} \\
\midrule
Boton\_marcha & Bool & \%I1.0 \\
Boton\_paro & Bool & \%I1.1 \\
Selector\_modo & Bool & \%I1.2 \\
Boton\_rearme & Bool & \%I1.3 \\
Sensor\_ind\_bien\_pos & Bool & \%I1.4 \\
Seta\_emergencia & Bool & \%I1.5 \\
Sensor\_palet\_izq & Bool & \%I1.6 \\
Sensor\_palet\_dch & Bool & \%I1.7 \\
\bottomrule
\end{tabular}

\medskip
\subsection*{\texttt{UDT\_Salidas\_Basico}}
Tag \texttt{Salidas\_Basico}, desde \texttt{\%Q1.0}.\par\smallskip
\begin{tabular}{>{\ttfamily\small}p{5.2cm}>{\ttfamily\small}p{4.2cm}>{\ttfamily\small}p{2.4cm}}
\toprule
\textnormal{\textbf{Miembro}} & \textnormal{\textbf{Tipo}} & \textnormal{\textbf{Dirección}} \\
\midrule
Led\_boton\_marcha & Bool & \%Q1.0 \\
Led\_boton\_rearme & Bool & \%Q1.1 \\
Led\_q1\_ciclo\_auto & Bool & \%Q1.2 \\
Led\_q2\_emergencia & Bool & \%Q1.3 \\
Cinta\_dch & Bool & \%Q1.4 \\
Cinta\_izq & Bool & \%Q1.5 \\
Cinta\_velocidad\_baja & Bool & \%Q1.6 \\
Bajar\_stopper & Bool & \%Q1.7 \\
\bottomrule
\end{tabular}

\medskip
\subsection*{\texttt{UDT\_Entradas\_Stopper}}
Tag \texttt{Entradas\_Stopper}, desde \texttt{\%I42.0}.\par\smallskip
\begin{tabular}{>{\ttfamily\small}p{5.2cm}>{\ttfamily\small}p{4.2cm}>{\ttfamily\small}p{2.4cm}}
\toprule
\textnormal{\textbf{Miembro}} & \textnormal{\textbf{Tipo}} & \textnormal{\textbf{Dirección}} \\
\midrule
Sensor\_ind\_0 & Bool & \%I42.0 \\
Sensor\_ind\_1 & Bool & \%I42.1 \\
Sensor\_ind\_2 & Bool & \%I42.2 \\
Sensor\_ind\_3 & Bool & \%I42.3 \\
Reservado\_E2\_4 & Bool & \%I42.4 \\
Reservado\_E2\_5 & Bool & \%I42.5 \\
Reservado\_E2\_6 & Bool & \%I42.6 \\
Sensor\_stopper\_abajo & Bool & \%I42.7 \\
\bottomrule
\end{tabular}

\medskip
\subsection*{\texttt{UDT\_Entradas\_RFID}}
Tag \texttt{Entradas\_RFID}, desde \texttt{\%I10.0}.\par\smallskip
\begin{tabular}{>{\ttfamily\small}p{5.2cm}>{\ttfamily\small}p{4.2cm}>{\ttfamily\small}p{2.4cm}}
\toprule
\textnormal{\textbf{Miembro}} & \textnormal{\textbf{Tipo}} & \textnormal{\textbf{Dirección}} \\
\midrule
Tag\_datos\_escritos & Bool & \%I10.0 \\
Tag\_datos\_leidos & Bool & \%I10.1 \\
Reservado\_E4\_2 & Bool & \%I10.2 \\
Reservado\_E4\_3 & Bool & \%I10.3 \\
Tag\_antena\_apagada & Bool & \%I10.4 \\
Tag\_detectado & Bool & \%I10.5 \\
Tag\_error & Bool & \%I10.6 \\
Tag\_proceso\_listo & Bool & \%I10.7 \\
Tag\_codigo\_error & Byte & \%IB11 \\
Tag\_direccion\_confirmada & Word & \%IW12 \\
Tag\_datos\_leidos\_escritos & Array[0..27] of Byte & \%I14.0 ... \%I41.7 \\
\bottomrule
\end{tabular}

\medskip
\subsection*{\texttt{UDT\_Salidas\_RFID}}
Tag \texttt{Salidas\_RFID}, desde \texttt{\%Q10.0}.\par\smallskip
\begin{tabular}{>{\ttfamily\small}p{5.2cm}>{\ttfamily\small}p{4.2cm}>{\ttfamily\small}p{2.4cm}}
\toprule
\textnormal{\textbf{Miembro}} & \textnormal{\textbf{Tipo}} & \textnormal{\textbf{Dirección}} \\
\midrule
Tag\_escribir & Bool & \%Q10.0 \\
Tag\_leer & Bool & \%Q10.1 \\
Tag\_reservado\_2 & Bool & \%Q10.2 \\
Tag\_reservado\_3 & Bool & \%Q10.3 \\
Tag\_antena\_apagar & Bool & \%Q10.4 \\
Tag\_reservado\_5 & Bool & \%Q10.5 \\
Tag\_reservado\_6 & Bool & \%Q10.6 \\
Tag\_reservado\_7 & Bool & \%Q10.7 \\
Reservado\_AB5 & Byte & \%QB11 \\
Tag\_direccion\_leer\_escribir & Word & \%QW12 \\
Tag\_datos\_a\_escribir & Array[0..27] of Byte & \%Q14.0 ... \%Q41.7 \\
\bottomrule
\end{tabular}

\medskip


\begin{ojo}
Estos UDT tienen que estar \textbf{bit a bit} como se muestran. Un miembro movido de sitio cambia el sensor al que se refiere. Por eso los tipos de la librería no se editan en los proyectos.
\end{ojo}
